% GeoMERIT — BDCC revision (merged final).
% Built on the official MDPI LaTeX template (Definitions/mdpi.cls).
% Text reorganized to match the method-paper narrative of 最新论文.pdf / v2.pdf.
% Section 3 figures updated: Figure 1 now combines the new well-location map,
% the long-tail class distribution panel, and the new family-grouped missingness matrix.
\documentclass[bdcc,article,submit,pdftex,moreauthors]{Definitions/mdpi}

\firstpage{1}
\makeatletter
\setcounter{page}{\@firstpage}
\makeatother
\pubvolume{1}
\issuenum{1}
\articlenumber{0}
\pubyear{2026}
\copyrightyear{2026}
\datereceived{ }
\daterevised{ }
\dateaccepted{ }
\datepublished{ }

\Title{GeoMERIT: Missingness-Aware and Penalty-Guided Robust Lithology Prediction from Well Logs}

\newcommand{\orcidauthorA}{0000-0000-0000-000X}
\Author{Anonymous Author $^{1,}$*\orcidA{}}
\AuthorNames{Anonymous Author}
\address{$^{1}$ \quad Affiliation 1; email@example.com}
\corres{Correspondence: email@example.com}

\abstract{Lithology prediction from field well logs must be performed with incomplete,
heterogeneous, and geologically constrained observations. Entire log curves may be unavailable
for a well, acquisition settings vary among wells, lithology labels are highly imbalanced,
formation-boundary samples are uncertain, and the severity of a mistake depends on the
geological relation between the predicted and true classes. These conditions can make a
conventional classifier perform well under random row-wise splits while generalizing poorly to
unseen wells, incomplete curve suites, rare lithologies, and boundary intervals. We present
GeoMERIT, a missingness-aware and penalty-guided method for robust lithology prediction.
GeoMERIT addresses the main sources of well-log uncertainty at four stages of the prediction
pipeline. Missing curves and anomalous responses are encoded as part of a learnable geological
state, together with curve-suite identity, within-well normalization, and local depth context,
rather than being reconstructed as ordinary numeric values. Class support, boundary status,
and interpreter confidence define sample weights for training. The resulting class
probabilities are refined through a three-backend gradient-boosted ensemble, a coarse-to-fine
geological hierarchy, and rare-class experts, reducing the dominance of head-class priors over
rare and geologically adjacent classes. The final label is selected by minimizing expected
geological cost under the FORCE penalty matrix rather than by posterior argmax. On the FORCE
2020 dataset (118 wells, 12 lithology classes, 1{,}431{,}383 valid labeled samples), under
strict 10-fold validation grouped by well, GeoMERIT reaches Weighted F1 $=0.7526$, Boundary F1
$=0.4567$, and Penalty $=-0.5886$, improving on the published RFE+XGBoost reference by
$+3.5\%$, $+11.4\%$, and $+6.4\%$ simultaneously, with a tail-class mean F1 of $0.5061$.
These concurrent gains show that coordinated treatment of data quality, class imbalance,
boundary uncertainty, posterior structure, and geological cost improves robustness under
realistic well-log conditions.}

\keyword{robust lithology prediction; well logs; missing data; long-tailed learning;
boundary samples; cost-sensitive decision; gradient-boosted decision trees; FORCE 2020}


% Center short figure captions under the corresponding figure without moving detailed explanations into captions.
\newcommand{\ShortCenteredFigCaption}[2]{%
  \refstepcounter{figure}%
  \label{#2}%
  \vspace{0.08cm}%
  \begin{center}
  \small\textbf{Figure~\thefigure.}~#1
  \end{center}%
}

\begin{document}
\sloppy

\section{Introduction}
\label{sec:intro}
Lithology prediction is a foundational task in reservoir characterization, geological
modeling, and hydrocarbon development evaluation~\cite{rider2011,ellis2007,doveton2014}. Each
well log records the formation response to a battery of physical measurements---gamma ray
(GR), bulk density (RHOB), neutron porosity (NPHI), photoelectric factor (PEF), sonic transit
time (DTC), resistivity (RDEP), and others---along depth, and machine learning methods are
now routinely used to map these curves to lithology labels at
scale~\cite{hall2016facies,merembayev2021,dev2019formation,carvalho2026benchmark}. Implicitly,
however, most of this work inherits ordinary supervised-classification assumptions, including
complete inputs, reasonably balanced classes, identically distributed training and test samples,
and equally costly errors. Real well-log data satisfy none of these
assumptions.

Field well-log datasets combine several coupled sources of uncertainty. Curve availability is
systematic and well dependent, reflecting tool strings, vendors, acquisition vintages, and
borehole conditions. Measured responses may also shift across wells because of calibration and
borehole-environment effects, and the records can contain anomalous or noisy intervals. At the
label level, lithology distributions are highly long-tailed, formation-boundary samples are
uncertain, and the geological severity of an error depends on which lithologies are confused.
The FORCE 2020 benchmark~\cite{bormann2020force} provides a public dataset in which these
conditions occur together. Among its 1{,}431{,}383 valid labeled samples from 118 Norwegian
North Sea wells, Shale accounts for 61.3\% of labels while Basement contributes 141 rows
(0.01\%). This head-to-tail ratio of roughly $6220{:}1$ coexists with entire log families
missing in specific wells (Figure~\ref{fig:missingness_matrix}), and the dataset includes a
$12{\times}12$ penalty matrix in which, for example, predicting Sandstone/Shale for true
Sandstone is a mild error whereas predicting Halite for the same rock is a severe one. The
modeling consequence is a generalization problem rather than merely a lower accuracy score. A
classifier can perform well under random row-wise splits, where well-specific acquisition
signatures and local stratigraphic continuity leak between training and test rows~\cite{rosenblatt2024,roberts2017cv,mila2022nndm},
yet degrade when applied to an unseen well, an incomplete curve suite, a rare lithology, or a
formation-boundary interval.

Prior work addresses important parts of this problem, but the components are usually treated
separately. Imputation-centered workflows~\cite{vanbuuren2011mice,stekhoven2012missforest,liu2025wellgan,xie2024cgan}
reconstruct plausible curve values; in well logs, however, the missingness pattern can itself
be informative because the decision to run a curve reflects vintage, vendor, borehole
condition, and operational choices that may correlate with geology~\cite{chen2026rfe,little2019stats,josse2024missing}.
Class-imbalance remedies re-weight or over-sample rare
classes~\cite{chawla2002smote,lin2017focal,cui2019classbalanced,he2024imbalanced}, but they
rarely distinguish formation-boundary samples, where labels are uncertain and interpretive
value is high. Deep architectures can learn rich representations, but they often require
complete curve suites and large labeled corpora, and they have not shown consistent superiority
over tree ensembles when evaluation is grouped by
well~\cite{wang2024dnnlitho,grinsztajn2022trees,shwartzziv2022tabular,mcelfresh2024benchmark}.
In addition, most published systems decode by posterior argmax, which treats all errors as
exchangeable even when an official cost matrix assigns different geological penalties~\cite{elkan2001cost,berger1985decision}.
A recent robust-feature-engineering (RFE) study with quality-aware XGBoost~\cite{chen2026rfe}
established the value of treating missingness as signal at the feature level, raising Weighted
F1 from $0.6640$ to $0.7272$ on FORCE 2020 under well-grouped validation. The remaining
question is whether robustness improves further when data quality, long-tail structure,
boundary uncertainty, and geological cost are modeled jointly across representation, learning,
posterior modeling, and decision making.

Motivated by this gap, we present GeoMERIT (Missingness-encoding, Ensembling,
Robust-features, Imbalance-aware, Tail-correction), a robust lithology prediction method for
incomplete, heterogeneous, and long-tailed well-log data. GeoMERIT distributes robustness
across the prediction pipeline, assigning a specific mechanism to each source of degradation
(Section~\ref{sec:method}, Figure~\ref{fig:framework}, Table~\ref{tab:layers}). Curve absence,
anomalous responses, curve-suite combinations, within-well normalized responses, and local
depth context are first encoded as an explicit geological state, so data quality is treated as
predictive information rather than as a preprocessing defect. On that representation, joint
class, boundary, and confidence weighting increases the influence of rare lithologies and
formation transitions relative to frequent interior intervals. The posteriors are then refined
through a three-backend tree ensemble, a coarse-to-fine geological hierarchy, and dedicated tail
experts, allowing common, adjacent, and rare lithologies to be modeled at different levels of
resolution. The final label is selected by minimizing expected geological cost under the FORCE
penalty matrix rather than by maximum-posterior classification. Throughout, well-grouped
validation is used as the \emph{test} of robustness: it simulates prediction on wells the model
has never seen, but it does not define the method itself. The objective is to remain stable
under the missingness, heterogeneity, long-tail structure, and cost asymmetry that any new well
may present. Under this protocol GeoMERIT improves on the RFE+XGBoost reference
simultaneously on Weighted F1 ($0.7272\!\to\!0.7526$, $+3.5\%$), Boundary F1
($0.4100\!\to\!0.4567$, $+11.4\%$), and FORCE Penalty ($-0.6289\!\to\!-0.5886$, $+6.4\%$),
while raising the tail-class mean F1 to $0.5061$. The framework is implemented entirely with
CPU-trainable tree ensembles, supporting reproducibility without specialized accelerator
hardware.

The overarching contribution is a single robust prediction framework that treats missingness,
acquisition heterogeneity, long-tailed supervision, boundary ambiguity, and asymmetric
geological cost as model inputs, learning constraints, posterior structure, and decision costs.
The framework rests on four testable design choices.
\begin{enumerate}
\item \textbf{A missingness-aware robust geologic state representation.} Rather than imputing
absent curves---which assumes missingness is ignorable---we encode curve absence, anomalous
responses, curve-suite identity, well-specific normalization, local depth context, and
stratigraphic metadata into a unified state vector, so the model can learn which curves are
missing, in which combinations, and how those patterns relate to lithology. The directly
testable consequence is graceful degradation on feature-deficient wells
(Sections~\ref{sec:repr} and \ref{sec:rq3}).
\item \textbf{A long-tail- and boundary-enhanced learning mechanism.} A multiplicative sample
weight combining effective-number class weighting, boundary up-weighting, and label-confidence
discounting increases the influence of rare lithologies and formation transitions relative to
frequent interior intervals, raising tail mean F1 to $0.5061$ (about $1.6\times$ the
single-model level) and contributing to the largest relative headline gain, $+11.4\%$
Boundary F1 (Sections~\ref{sec:learning} and \ref{sec:rq4}).
\item \textbf{Structured posterior lithology modeling.} Rather than relying only on a single
12-class posterior, GeoMERIT fuses LightGBM, XGBoost, and CatBoost posteriors geometrically,
mixes in a coarse-to-fine geological hierarchy over six super-classes, and gates six
one-vs-rest tail experts, so rare-class posterior mass is less dominated by head-class priors
and geologically adjacent classes are handled by dedicated experts (Sections~\ref{sec:posterior}
and \ref{sec:rq2}).
\item \textbf{A geological-cost-guided decision layer.} Lithology prediction is converted from
maximum-probability classification to minimum-geological-risk decision. Logit adjustment
corrects the long-tail prior shift, while Bayes-risk decoding over the official FORCE penalty
matrix selects the label with the lowest expected geological cost. To our knowledge GeoMERIT
is the first FORCE 2020 system whose headline predictions are produced this way, and the
$+6.4\%$ Penalty gain quantifies the effect (Sections~\ref{sec:decision} and \ref{sec:rq5}).
\end{enumerate}

The remainder of the paper is organized as follows. Section~\ref{sec:related} positions
GeoMERIT against four families of prior work. Section~\ref{sec:data} describes the dataset,
the metrics, and the well-grouped robustness protocol. Section~\ref{sec:method} formalizes
the four framework layers. Section~\ref{sec:setup} details the experimental setup and
comparison fairness. Section~\ref{sec:results} reports results, moving from overall
held-out-well performance through the sources of the gains to behavior on feature-deficient
wells, rare lithologies, and the decision layer. Section~\ref{sec:discussion} analyzes failure
modes and limitations, and
Section~\ref{sec:conclusion} concludes.

\section{Related Work}
\label{sec:related}
We group prior work into four families, each defined by the aspect of real-data complexity
it addresses, and locate GeoMERIT by what each family leaves open.

\subsection{Well-Log Lithology Classification and Benchmarking}
Machine learning lithology classification has progressed from early facies
challenges~\cite{hall2016facies} through systematic algorithm
comparisons~\cite{merembayev2021,dev2019formation} to deep architectures including CNNs,
recurrent models, transformers, diffusion-based augmentation, and multi-scale hierarchical
networks~\cite{wang2024dnnlitho,chen2025geological,zhao2025diffusion,zhan2025mchgeo}.
Interpretability-oriented hybrids such as SHAP-based meta-learning
transformers~\cite{sun2024shapmeta,lundberg2017shap} and refined-lithology
pipelines~\cite{shi2023refined} demonstrate strong within-dataset accuracy. Benchmark studies
make the comparability problem explicit. Published scores are frequently incomparable because
splits, label sets, and metrics differ, and reported performance can swing by 10--20 absolute
F1 points with the validation protocol alone, because the unit of geological generalization
is the well, not the row~\cite{carvalho2026benchmark,rosenblatt2024}. What remains open is a
method \emph{designed}---not merely evaluated---for stability under the data conditions that
cause those swings. GeoMERIT adopts the benchmark literature's protocol discipline and builds
the corresponding robustness mechanisms into every stage of the model.

\subsection{Missingness and Data-Quality Modeling in Well Logs}
The dominant response to incomplete logs is reconstruction, most commonly through chained-equation
and forest imputation~\cite{vanbuuren2011mice,stekhoven2012missforest} or through generative
imputation with sequence GANs and conditional GANs~\cite{liu2025wellgan,xie2024cgan}. Reconstruction is
appropriate when missingness is ignorable. Well-log missingness, however, is structurally
non-random because whole curves are absent as a result of tool-string choices that correlate
with vintage, operator, and target formation~\cite{chen2026rfe}. Statistical learning theory shows
that supervised learners can remain consistent when missingness is encoded as input rather
than replaced by imputed values~\cite{little2019stats,josse2024missing}, and the RFE
study~\cite{chen2026rfe} demonstrated the practical value of this position on FORCE 2020.
The open question is how far the ``missingness as signal'' paradigm extends beyond features
into training and decision-making. GeoMERIT extends it through representation, weighting,
posterior structure, and decoding.

\subsection{Long-Tailed and Boundary-Aware Learning}
General-purpose remedies for class imbalance include synthetic
oversampling~\cite{chawla2002smote}, focal loss~\cite{lin2017focal}, margin-based
losses~\cite{cao2019ldam}, effective-number class weighting~\cite{cui2019classbalanced}, and
prior-corrected logit adjustment~\cite{menon2021logit}; surveys document their
trade-offs~\cite{he2024imbalanced}. Two domain constraints restrict their transfer to well
logs. First, synthetic interpolation of log rows can violate depth continuity and
stratigraphic consistency, generating physically implausible samples; we therefore avoid
synthetic oversampling entirely. Second, lithology boundary intervals contain both label
ambiguity and practical interpretive value, so imbalance handling must not sacrifice boundary
fidelity---a requirement most imbalance methods rarely examine. GeoMERIT combines
effective-number weighting~\cite{cui2019classbalanced} with explicit boundary up-weighting and
treats logit adjustment~\cite{menon2021logit} as a prior-shift correction applied at decision
time, where its temperature can be tuned against the geological cost directly.

\subsection{Cost-Sensitive and Geology-Aware Decision Making}
Cost-sensitive learning theory establishes that when misclassification costs are asymmetric,
the optimal decision minimizes expected cost rather than maximizing posterior
probability~\cite{elkan2001cost,berger1985decision,andrade2019bayes,zadrozny2002multiclass}.
FORCE 2020 operationalizes geological cost in its penalty matrix, yet most published
systems---including the RFE reference~\cite{chen2026rfe}---still decode by argmax and report
the penalty only as an evaluation metric. Sequence decoders such as Viterbi smoothing and
(semi-)Markov CRFs~\cite{viterbi1967,lafferty2001crf,sarawagi2004semimarkov} optimize label
continuity but not geological cost. To our knowledge, GeoMERIT is the first FORCE 2020 system
whose headline predictions are produced by Bayes-risk decoding over the official penalty
matrix, coupled with logit adjustment in a single, jointly tuned decision layer.

\section{Dataset and Evaluation Protocol}
\label{sec:data}

\subsection{Dataset}
\label{sec:dataset}
We use the public FORCE 2020 lithofacies dataset~\cite{bormann2020force}, which contains
$1{,}431{,}383$ valid labeled samples from 118 wells and 12 lithology classes. Each sample is
indexed by measured depth and may contain up to 20 well-log curves, including GR, SGR, RHOB,
DRHO, NPHI, PEF, DTC, DTS, RDEP, RSHA, RMED, RXO, RMIC, SP, CALI, BS, DCAL, ROP, ROPA, and
MUDWEIGHT. We also use well-head coordinates and NPD stratigraphic tables, including GROUP,
FORMATION, and casing-depth information. These variables define a realistic lithology
prediction setting in which the model must infer rock type from heterogeneous curve suites,
uneven sample support, and geological labels whose errors have unequal consequences.

The first source of heterogeneity is spatial and operational. The 118 wells are distributed
across multiple North Sea blocks, and the available curve suites differ across those blocks.
Figure~\ref{fig:well_locations} therefore serves as the entry point to the dataset. It shows
where the wells are located, how many labeled samples each well contributes, and how much of
its log suite is missing on average.

\begin{figure}[H]
\centering
\includegraphics[width=0.78\linewidth]{figures/fig01_well_locations_v3.pdf}
\ShortCenteredFigCaption{Spatial distribution of FORCE 2020 wells.}{fig:well_locations}
\end{figure}

In Figure~\ref{fig:well_locations}, marker color represents the per-well mean missing rate,
and marker size represents the number of labeled samples in each well. The figure shows that
data quality is not uniformly distributed across the basin. Nearby wells can share similar
missingness patterns, while different blocks exhibit different acquisition regimes. This
means that a held-out well is not merely another independent row sample; it may correspond to
a new logging regime with a different combination of available and unavailable curves.

The second difficulty is the supervision structure. Even when curves are present, the labels
are highly imbalanced, so a model trained only to maximize aggregate accuracy can be dominated
by the most frequent lithologies. Figure~\ref{fig:lithology_longtail} isolates this label-side
problem before we examine the input-side missingness in more detail.

\begin{figure}[H]
\centering
\includegraphics[width=0.82\linewidth]{figures/fig02_lithology_distribution_line.pdf}
\ShortCenteredFigCaption{Long-tailed lithology class distribution.}{fig:lithology_longtail}
\end{figure}

Figure~\ref{fig:lithology_longtail} shows a long-tailed distribution spanning roughly four
orders of magnitude. Shale contains $877{,}043$ samples, accounting for 61.3\% of all valid
labels, whereas Basement contains only 141 samples, or about 0.01\%. The resulting
head-to-tail ratio is approximately $6220{:}1$, and the distribution has a Gini coefficient
of about 0.77. This imbalance explains why ordinary classifiers tend to over-optimize thick,
head-class intervals while under-learning rare lithologies and thin stratigraphic events.

The spatial heterogeneity in Figure~\ref{fig:well_locations} and the label imbalance in
Figure~\ref{fig:lithology_longtail} interact with a third property. The log curves are not
missing independently. Figure~\ref{fig:missingness_matrix} makes this input-side structure
explicit by arranging wells and curves as a missing-fraction matrix.

\begin{figure}[H]
\centering
\includegraphics[width=0.86\linewidth]{figures/fig03_missingness_vAB_merged.pdf}
\ShortCenteredFigCaption{Well-by-curve missing-fraction matrix.}{fig:missingness_matrix}
\end{figure}

Figure~\ref{fig:missingness_matrix} groups the curves into the petrophysical families Gamma,
Density/Neutron, Sonic, Resistivity, SP, Caliper, and Drilling/mud. Wells are sorted by their
mean missing rate, while the marginal bars summarize per-curve and per-well missingness. The
block patterns show that missingness is structured by both well and curve family. Entire log
families, such as Sonic or Resistivity, can be absent in specific wells, whereas other curve
families remain available. This is the key reason we do not treat missing values as random
noise to be imputed away. In this dataset, missingness carries information about tool-string
selection, acquisition vintage, and operational conditions, all of which can correlate with
geology. The representation layer in Section~\ref{sec:repr} is therefore designed to encode
missingness and curve-suite identity directly rather than erase them during preprocessing.

Finally, robust lithology prediction also requires a geology-aware evaluation target. The
FORCE 2020 benchmark provides an official $12{\times}12$ penalty matrix, shown in
Figure~\ref{fig:penalty}. This matrix assigns small costs to sedimentologically adjacent
confusions, such as Sandstone versus Sandstone/Shale, and larger costs to distant confusions,
such as Sandstone versus Halite.

\begin{figure}[H]
\centering
\includegraphics[width=0.72\linewidth]{figures/fig2_penalty_matrix.pdf}
\ShortCenteredFigCaption{Geological penalty matrix.}{fig:penalty}
\end{figure}

The penalty matrix in Figure~\ref{fig:penalty} makes the evaluation problem different from
ordinary multi-class classification. Entry $A_{ij}$ is the cost of predicting class $j$ when
the true class is $i$; the diagonal cost is zero, within-environment confusions are relatively
cheap, and cross-environment confusions can cost up to 4. Consequently, posterior argmax is
not necessarily the best geological decision. GeoMERIT uses this matrix at the decision layer
by minimizing expected geological penalty rather than treating all errors as equal.

\subsection{Evaluation Metrics: Three Facets of Robustness}
\label{sec:metrics}
We report five complementary metrics. Weighted F1 measures aggregate classification quality,
whereas Macro F1 gives the unweighted mean over the 12 classes. Boundary F1 is computed only
on samples within one depth step of a lithology change and therefore measures stratigraphic
transition quality. Penalty is the mean of $-A_{y\hat y}$ under the FORCE matrix, with values
closer to zero indicating less severe geological errors. Tail mean F1 is the unweighted mean
F1 over the six rarest classes, namely Tuff, Halite, Coal, Dolomite, Anhydrite, and Basement. The three headline metrics probe three distinct facets
of robust prediction and are not interchangeable. Weighted F1 reflects overall prediction
quality and can rise simply by getting more Shale rows right. Boundary F1 reflects stability
in formation-transition zones and rises only by resolving transitions. Penalty reflects the
geological severity of the errors that remain and improves only by avoiding geologically
distant confusions. A method that improves one facet by sacrificing another is not robust;
we therefore evaluate all three jointly.

\subsection{Held-Out-Well Validation as a Robustness Test}
\label{sec:protocol}
We evaluate exclusively under 10-fold GroupKFold cross-validation with well identity as the
grouping key. This protocol is the experimental condition that tests robustness, not part of
the method's definition. GeoMERIT's goal is not to adapt one specific well to another, but to
remain stable under the curve missingness, acquisition heterogeneity, long-tail structure,
and cost asymmetry that any previously unseen well may present, and holding out whole wells
is how that stability is measured. Random row-wise splits place adjacent depth samples from
the same well on both sides of the train/test boundary, allowing models to exploit
well-specific acquisition signatures and local stratigraphic continuity; such splits estimate
interpolation performance within already-logged wells and systematically overstate practical
performance~\cite{rosenblatt2024,roberts2017cv,mila2022nndm}. All feature construction that
uses depth context---window statistics, gradients, boundary masks---is computed strictly
within each well sequence, so no information crosses well boundaries at any stage.

This protocol deliberately exposes a rare-class limitation. For classes concentrated in a few
wells, some training folds contain no positive examples. Basement (141 rows) is the extreme
case. This is not an artifact of the split but a property of the prediction problem itself---a
new well may contain a lithology the training fleet never sampled---so we report Basement
separately and discuss open-set treatment as the appropriate future remedy
(Section~\ref{sec:failures}).

\subsection{Statistical Reporting}
\label{sec:stats}
Headline values are pooled out-of-fold (OOF) metrics over all $1{,}431{,}383$ predictions.
Because pooled values hide fold-to-fold variation, we additionally report the mean
$\pm$ standard deviation across the 10 folds and, for each headline metric, a one-sided
one-sample Wilcoxon signed-rank test~\cite{wilcoxon1945} of the 10 fold values against the
published RFE+XGBoost reference value (alternative: GeoMERIT exceeds the reference). This
test is conservative in one respect (the reference is a pooled point estimate, not a matched
per-fold series) and we interpret it accordingly in Section~\ref{sec:rq1}.

\section{GeoMERIT: A Four-Layer Robust Prediction Framework}
\label{sec:method}
GeoMERIT is organized as a stage-matched robust prediction framework. The main sources of
degradation in field well-log data---\emph{systematic curve missingness, acquisition
heterogeneity, long-tailed supervision, boundary ambiguity, and asymmetric geological
cost}---are assigned to the pipeline stages where they can be modeled directly. The framework
therefore comprises four layers (Table~\ref{tab:layers}, Figure~\ref{fig:framework}). Let $x_i$
denote the log response vector at depth sample $i$, $w(i)$ its training weight,
$p(y\,|\,\phi(x_i))$ the class posterior over the state representation $\phi(x_i)$, and $A$
the FORCE penalty matrix. GeoMERIT first constructs a missingness-aware state vector
$\phi(x_i)$ that preserves data-quality and curve-suite information
(Section~\ref{sec:repr}). It then estimates posteriors under class-, boundary-, and
confidence-aware weighting (Section~\ref{sec:learning}), refines them with a model ensemble, a
geological hierarchy, and tail experts (Section~\ref{sec:posterior}), and finally predicts
\begin{equation}
\hat y_i = \arg\min_{j\in\{1,\ldots,12\}} \sum_{k=1}^{12} A_{kj}\;\hat p\!\bigl(y=k\,\big|\,\phi(x_i)\bigr),
\label{eq:bayes}
\end{equation}
after long-tail logit adjustment (Section~\ref{sec:decision}). Equation~(\ref{eq:bayes}) is
the classical Bayes decision rule under an asymmetric loss~\cite{berger1985decision}. It
reduces to argmax when $A$ is the 0--1 loss and departs from argmax exactly where
geological costs are asymmetric. This design differs from a conventional feature-engineering
workflow because robustness is distributed across representation, learning, posterior modeling,
and decision making rather than confined to the feature table.

\begin{table}[H]
\caption{The four-layer GeoMERIT framework. Each layer addresses one condition of field
well-log data that destabilizes conventional lithology classifiers.\label{tab:layers}}
\begin{tabularx}{\textwidth}{lXXl}
\toprule
\textbf{Layer} & \textbf{Field-data condition} & \textbf{GeoMERIT mechanism} & \textbf{Section}\\
\midrule
Representation & Systematic curve missingness, anomalous responses, inter-well acquisition shift &
Sentinel encoding, missing/outlier indicators, within-well robust normalization, multi-scale
depth context, curve-suite fingerprint & \ref{sec:repr}\\
Learning & Long-tailed classes, ambiguous boundary samples, uncertain labels &
Multiplicative class $\times$ boundary $\times$ confidence sample weighting & \ref{sec:learning}\\
Posterior & Single multiclass posterior can underresolve rare and adjacent classes &
Three-backend GBDT geometric fusion, coarse-to-fine geological hierarchy, one-vs-rest tail experts &
\ref{sec:posterior}\\
Decision & Asymmetric geological misclassification cost ignored by argmax &
Logit adjustment plus Bayes-risk decoding over the FORCE penalty matrix & \ref{sec:decision}\\
\bottomrule
\end{tabularx}
\end{table}

\begin{figure}[H]
\centering
\includegraphics[width=\linewidth]{figures/geo_v2.pdf}
\caption{The GeoMERIT four-layer robust design. Each real-data challenge (left) is addressed
by one mechanism layer (right). The representation layer encodes systematic curve-suite
missingness and anomalous responses, the learning layer reduces the dominance of head-class
and interior samples, the posterior layer improves resolution for rare and adjacent classes
through a three-backend ensemble, a coarse-to-fine geological hierarchy, and tail experts, and
the decision layer replaces argmax with expected-cost minimization over the FORCE penalty
matrix. Robustness is then evaluated on unseen wells under strict 10-fold GroupKFold by well.}
\label{fig:framework}
\end{figure}

\subsection{Missingness-Aware Geologic State Representation (Representation Layer)}
\label{sec:repr}
In field well-log collections, curve availability is often systematic rather than random. The
curves recorded for a well reflect the tool string, acquisition vintage, vendor practice, and
borehole conditions; as a result, entire log families may be unavailable across a well rather
than sporadically missing at individual depths (Figure~\ref{fig:missingness_matrix}). Treating
this absence only as a preprocessing gap to be filled by mean or regression imputation has two
limitations. The imputed values can introduce well-correlated error, and the acquisition
information carried by the missingness pattern is removed. Because that pattern helps
characterize the logging regime of a new well, GeoMERIT encodes missingness directly as part of
the learnable geological state.

Concretely, for every curve $c$ we detect missing entries (NaN/Inf) and anomalies (within-well
robust $z$-score $|z|>5$) and produce a sentinel-encoded copy $c_\text{enc}$ in which missing
values become $-999$ and outliers $-777$, alongside binary indicators $c_\text{is\_missing}$
and $c_\text{is\_outlier}$; per-row aggregates of the missing count, missing ratio, outlier
count, and outlier ratio summarize the data quality at each depth. In a GBDT, sentinel values
can be routed through dedicated split branches, allowing rules such as ``if PEF is absent,
use GR and NPHI instead'' to be learned without explicit rule engineering. The model therefore
uses both the observed log responses and the curve-availability pattern.

Within each well $w$ and curve $c$ we add a robust normalization
\begin{equation}
z_{w,c}(i) = \frac{c(i)-\operatorname{median}_w(c)}{\operatorname{MAD}_w(c)+\epsilon},\qquad \epsilon=10^{-9},
\label{eq:zscore}
\end{equation}
with a standard-deviation fallback when $\operatorname{MAD}_w(c)=0$. Both the absolute
encoded value and the well-relative value are retained. The absolute response separates
lithologies with distinctive global signatures (e.g., Halite), while the well-relative offset
compensates for calibration and borehole-environment drift between wells---precisely the
inter-well heterogeneity that destabilizes transfer to a new well.

Local depth context is captured by window statistics and stabilized gradients. For the eight
primary curves (GR, RHOB, NPHI, PEF, DTC, RDEP, RMED, SP) and radii $r\in\{2,5,10\}$ we
compute window mean, standard deviation, minimum, maximum, and endpoint difference
($8\!\times\!3\!\times\!5=120$ features), plus the depth-normalized gradient
\begin{equation}
\nabla c(i) = \frac{c(i)-c(i-1)}{\text{DEPTH}(i)-\text{DEPTH}(i-1)+\epsilon}.
\label{eq:grad}
\end{equation}
All context features are computed strictly within each well. Four diagnostic combinations
from the classical interpretation playbook (density--porosity separation $\text{RHOB}-\text{NPHI}$;
the ratios $\text{DTC}/|\text{RHOB}|$, $\text{GR}/\log(1+\text{RDEP})$, and
$\text{PEF}/|\text{RHOB}|$) inject petrophysical priors for separating carbonates from
clastics~\cite{rider2011,doveton2014}.

Finally, the per-well curve-availability bit string is hashed into a categorical
\texttt{curve\_suite\_id}---a single feature that captures vendor-, era-, and tool-driven
variation---and well-head coordinates, distance to the fleet centroid, NPD GROUP and
FORMATION labels, and casing-depth differences complete the state vector $\phi(x_i)$. Missing
stratigraphic labels are encoded as the explicit category ``Unknown'' rather than dropped.

\subsection{Long-Tail- and Boundary-Enhanced Risk Weighting (Learning Layer)}
\label{sec:learning}
Under an unweighted objective, the loss is dominated by high-support interior classes. Shale
alone supplies $61\%$ of the samples, so equal weighting allows the optimizer to reduce loss
primarily through the majority lithology, while boundary intervals and tail classes contribute
comparatively little to the objective. The learning layer therefore assigns each training
sample a multiplicative weight
\begin{equation}
w(i) = w_\text{class}\!\bigl(y(i)\bigr)\cdot w_\text{bnd}(i)\cdot w_\text{conf}(i).
\label{eq:weight}
\end{equation}
The class term uses the effective-number-of-samples
formulation~\cite{cui2019classbalanced},
\begin{equation}
w_\text{class}(k) = \min\!\left(\frac{1-\rho}{1-\rho^{n_k}},\; w_\text{max}\right),\qquad \rho=0.9995,\quad w_\text{max}=8.0,
\label{eq:effn}
\end{equation}
where $n_k$ is the class support in the training fold; weights are rescaled to unit mean.
Equation~(\ref{eq:effn}) functions as a saturation-aware inverse frequency. Once a class has
enough samples that new ones are redundant, additional support stops reducing its weight,
while the cap $w_\text{max}$ prevents the rarest classes from destabilizing training.
Boundary samples (label differs from an immediate depth neighbor) receive weight $1.2$ and
interior samples $0.7$, shifting a fixed attention budget toward transitions without
sacrificing the head. The confidence term maps the FORCE interpreter-confidence field
$1/2/3$ to $1.0/0.8/0.6$, discounting samples the original labelers themselves marked as
uncertain. We deliberately avoid synthetic oversampling~\cite{chawla2002smote} because interpolated
log rows can violate depth continuity and produce physically impossible curve combinations.

\subsection{Structured Posterior Modeling: Ensemble, Hierarchy, and Tail Experts (Posterior Layer)}
\label{sec:posterior}
A single 12-class posterior places all lithology labels in the same competition, although the
label space contains several discrimination regimes: separating geologically adjacent classes
(Limestone vs.\ Chalk), separating distant classes (Sandstone vs.\ Halite), and detecting an
extremely rare class such as Basement (141 samples). Under this shared normalization,
head-class priors can suppress posterior mass for rare classes. The posterior layer therefore
retains the multiclass classifier but augments it with three coordinated mechanisms that
increase resolution for rare and geologically adjacent lithologies.

The first mechanism replaces the lone classifier with three complementary backbones trained
per fold---LightGBM~\cite{ke2017lightgbm} on the full training partition ($\approx\!1.28$M rows),
XGBoost~\cite{chen2016xgboost} on a 750{,}000-row subsample, and
CatBoost~\cite{prokhorenkova2018catboost} on a 300{,}000-row subsample, both subsamples
preserving all tail rows and CatBoost run with $\texttt{classes\_count}=12$ so no class is
silently dropped from folds lacking it. Their inductive biases differ (leaf-wise growth with
native categorical support; depth-wise regularized histogram splitting; ordered boosting with
target statistics). Their posteriors are fused geometrically in log space as
\begin{equation}
\hat p(y=k\,|\,\phi)\propto\exp\!\left(\sum_m \alpha_m\log p^{(m)}(y=k\,|\,\phi)\right),\qquad
\boldsymbol\alpha=(1.0,\,0.8,\,0.8).
\label{eq:geomean}
\end{equation}
The geometric mean acts as a product of experts in which a class receives high fused probability only
if no backend strongly vetoes it, which is more robust than an arithmetic mean when one
backend is miscalibrated on a rare class.

The second mechanism introduces a geological hierarchy in the label space. The 12 classes are
mapped to six super-classes that mirror interpreter taxonomy (Table~\ref{tab:superclass}); a
coarse classifier predicts the super-class, per-group experts predict the fine class within it,
and the hierarchical posterior
\begin{equation}
p_\text{c2f}(y=k\,|\,\phi) = p\!\bigl(\text{super}(k)\,|\,\phi\bigr)\cdot p\!\bigl(k\,|\,\phi,\text{super}(k)\bigr)
\label{eq:c2f}
\end{equation}
is mixed with the flat posterior using weight $\mu=0.20$ (validation grid). The hierarchy
encourages competition within sedimentologically coherent groups. For example, Chalk and
Limestone are evaluated within the carbonate group, where carbonate-specific evidence can be
used without competing directly against the dominant Shale prior.

\begin{table}[H]
\caption{Geological super-classes used by the coarse-to-fine cascade.\label{tab:superclass}}
\begin{tabularx}{\textwidth}{lX}
\toprule
\textbf{Super-class} & \textbf{Fine classes}\\
\midrule
Clastic & Sandstone, Sandstone/Shale, Shale\\
Carbonate & Limestone, Chalk, Dolomite, Marl\\
Evaporite & Halite, Anhydrite\\
Coal & Coal\\
Basement & Basement\\
Tuff & Tuff\\
\bottomrule
\end{tabularx}
\end{table}

The third mechanism adds dedicated rare-class discriminators that work outside the 12-way
competition. For each tail class $k\in\mathcal T=\{\text{Tuff, Halite,
Coal, Dolomite, Anhydrite, Basement}\}$ a one-vs-rest LightGBM expert with $5\times$
up-weighted positives produces a score $s_k$. When the expert is confident, the corresponding
posterior entry is gated upward:
\begin{equation}
\tilde p(y=k\,|\,\phi) = \hat p(y=k\,|\,\phi)\!\left(1+\gamma_k\,\mathbf 1\bigl[s_k>\theta_k\bigr]\right),
\label{eq:gate}
\end{equation}
followed by renormalization. The gate restores posterior mass that the long tail suppressed,
but only where a dedicated binary discriminator---trained without the 12-way
competition---finds strong evidence. A single global pair $(\gamma,\theta)$ serves all tail
classes, with a separate, stricter Dolomite pair $(\gamma_\text{dol},\theta_\text{dol})$
because Dolomite trades off directly against the penalty (Section~\ref{sec:rq5}).

\subsection{Geological-Cost-Guided Decision Layer (Decision Layer)}
\label{sec:decision}
Two issues remain after posterior modeling. Argmax decoding treats every error as equally
costly, which is inconsistent with the asymmetric FORCE penalty matrix, and a posterior trained
on a long-tailed fold still carries the head-class prior into the final decision. The decision
layer addresses both issues before assigning a label. It first applies logit
adjustment~\cite{menon2021logit},
\begin{equation}
\hat p_\text{LA}(y=k\,|\,\phi)\propto \tilde p(y=k\,|\,\phi)\cdot \pi_k^{-\tau},
\label{eq:logitadj}
\end{equation}
where $\pi_k$ is the training-fold prior and $\tau\geq 0$ a temperature; this is a post-hoc
correction of the label prior shift between the long-tailed training distribution and the
evaluation objective, not a heuristic boost. The adjusted posterior then enters the
Bayes-risk rule of Equation~(\ref{eq:bayes}), so the final label is the one with the lowest
expected geological cost rather than the highest probability. The full decision layer has five
interpretable parameters---$\tau$, the global gate $(\gamma,\theta)$, and the Dolomite gate
$(\gamma_\text{dol},\theta_\text{dol})$---searched on a small validation grid (192
configurations), yielding $\tau=0.2$, $\gamma=1.2$, $\theta=0.35$,
$\gamma_\text{dol}=8.0$, $\theta_\text{dol}=0.05$. Each parameter has a direct operational
interpretation; for example, $\theta_\text{dol}=0.05$ with $\gamma_\text{dol}=8.0$ applies the
Dolomite expert only when its score exceeds the threshold, then gives that evidence strong
posterior weight. Domain experts can therefore retune the operating point without retraining.

An optional Viterbi pass over each well's depth sequence~\cite{viterbi1967,lafferty2001crf,sarawagi2004semimarkov}
can smooth the decoded labels with an empirical transition matrix; it is examined in
ablation but excluded from headline numbers to keep the decision layer purely pointwise and
tree-only.

\section{Experimental Setup}
\label{sec:setup}

\subsection{Compared Methods and Fairness}
\label{sec:compared}
The primary references are Baseline XGBoost (no feature engineering, no weighting),
RFE+XGBoost, and the RFE+\{Random Forest, CatBoost, LightGBM\} variants, all as published
in~\cite{chen2026rfe} and all evaluated on FORCE 2020 under a comparable well-grouped 10-fold
protocol, which makes them the strongest directly comparable tree baselines available. We use
the operating points reported by the original authors for these published baselines. GeoMERIT's
decision-layer grid is searched on validation folds under the same GroupKFold protocol used for
reporting, and comparisons are interpreted with this distinction in mind.

\subsection{Implementation and Scale}
\label{sec:impl}
All experiments ran on a single 112-core CPU machine with 251~GB RAM (no GPU). The software
stack is LightGBM, XGBoost, CatBoost, scikit-learn~\cite{pedregosa2011sklearn}, NumPy,
Pandas, and PyArrow. Each fold trains the three backends, the coarse-to-fine cascade, and six
tail experts; full 10-fold training takes about 3~h and the decoding-parameter search about
1~h. The feature table occupies $\approx\!12$~GB in Parquet form and peak training memory is
$\approx\!90$~GB.

\section{Results}
\label{sec:results}
We report the results in five steps. First, we evaluate whether GeoMERIT improves overall
classification, boundary recovery, and geological cost under held-out-well validation, and
whether these gains exceed fold-to-fold variability (Section~\ref{sec:rq1}). Second, we
examine the sources of the gains through published ablation rungs, backbone comparisons, and a
same-stack decision-layer contrast (Section~\ref{sec:rq2}). Third, we stratify test wells by
missingness to assess stability when curve suites are incomplete (Section~\ref{sec:rq3}).
Fourth, we report per-lithology behavior to identify rare-class and boundary limitations
(Section~\ref{sec:rq4}). Finally, we sweep the decision-layer parameters to assess whether the
cost-guided operating point is stable and interpretable (Section~\ref{sec:rq5}).

\subsection{Main Comparison Under Held-Out-Well Validation}
\label{sec:rq1}
The main evaluation asks whether the four-layer design improves three robustness indicators
under the same held-out-well protocol. We compare GeoMERIT with the published RFE+XGBoost
reference under 10-fold GroupKFold by well, taking Weighted F1 as overall classification
quality, Boundary F1 as stability in formation-transition zones, and Penalty as the geological
severity of the remaining errors. Under this protocol GeoMERIT improves all three metrics,
raising Weighted F1 from $0.7272$ to $0.7526$ ($+0.0254$; $+3.5\%$), Boundary F1 from
$0.4100$ to $0.4567$ ($+0.0467$; $+11.4\%$), and Penalty from $-0.6289$ to $-0.5886$
($+0.0403$; $+6.4\%$) (Table~\ref{tab:main}, Figure~\ref{fig:main}). The joint direction of
these changes is important because the metrics can trade off against one another. A model that
leans further toward the majority class can increase Weighted F1 while reducing boundary and
tail performance; a model that over-corrects rare classes can improve recall while increasing
geologically severe errors. The simultaneous gains indicate that GeoMERIT shifts the operating
point toward higher overall quality, better boundary stability, and lower geological cost, while
the tail mean F1 increases from roughly $0.31$ at the single-model baseline level to $0.5061$.

\begin{table}[H]
\caption{Main comparison under 10-fold GroupKFold by well. Reference values from~\cite{chen2026rfe};
GeoMERIT values are pooled out-of-fold over all $1{,}431{,}383$ samples. The third row
(``base ensemble $+$ Bayes'') uses the representation, learning, and posterior layers of
GeoMERIT but decodes with plain Bayes risk (no logit adjustment, no tail gates); contrasting
it with the full row isolates the decision layer in a same-stack ablation. $\Delta$ rows give
the change relative to RFE+XGBoost.\label{tab:main}}
\begin{tabularx}{\textwidth}{Xccccc}
\toprule
\textbf{Method} & \textbf{Weighted F1} & \textbf{Macro F1} & \textbf{Boundary F1} & \textbf{Penalty} & \textbf{Tail F1}\\
\midrule
Baseline XGBoost~\cite{chen2026rfe} & 0.6640 & -- & 0.3510 & $-0.7800$ & --\\
RFE + XGBoost~\cite{chen2026rfe} & 0.7272 & -- & 0.4100 & $-0.6289$ & --\\
GeoMERIT, base ensemble + Bayes decoding & 0.7446 & 0.5038 & 0.4423 & $-0.5885$ & 0.4737\\
\textbf{GeoMERIT, full decision layer} & \textbf{0.7526} & \textbf{0.5315} & \textbf{0.4567} & $\mathbf{-0.5886}$ & \textbf{0.5061}\\
\midrule
$\Delta$ absolute vs.\ RFE+XGBoost & $+0.0254$ & -- & $+0.0467$ & $+0.0403$ & --\\
$\Delta$ relative vs.\ RFE+XGBoost & $+3.5\%$ & -- & $+11.4\%$ & $+6.4\%$ & --\\
\bottomrule
\end{tabularx}
\end{table}

\begin{figure}[H]
\centering
\includegraphics[width=\linewidth]{figures/fig4_main_results.pdf}
\caption{GeoMERIT improves aggregate, boundary-sensitive, and cost-sensitive performance
simultaneously under held-out-well validation. Panels show Weighted F1, Boundary F1, and
Penalty for Baseline XGBoost, RFE+XGBoost~\cite{chen2026rfe}, and GeoMERIT; annotations give
the gain over the RFE reference. The Penalty axis is signed so that closer to zero is
better.}
\label{fig:main}
\end{figure}

Because pooled metrics can be influenced by fold composition, we also examine fold-level
dispersion and test each headline metric against the reference (Table~\ref{tab:perfold}).
Boundary F1 exceeds the reference in 9 of 10 folds (one-sided Wilcoxon signed-rank test,
$p=0.002$), so the boundary gain is individually significant. Weighted F1 exceeds it in 7 of
10 folds ($p=0.080$) and Penalty in 6 of 10 ($p=0.138$); these two gains are directionally
consistent but, taken alone, cannot be separated from fold heterogeneity at $\alpha=0.05$.
This is expected because fold composition varies geologically and the fold standard deviations
of $0.048$ and $0.094$ are larger than the corresponding increments. We therefore base the
headline interpretation on the joint pattern: simultaneous pooled improvement on all three
robustness facets, with the boundary component individually significant. Figure~\ref{fig:perfold}
shows the per-fold values behind this interpretation.

\begin{table}[H]
\caption{Fold-level dispersion of GeoMERIT (mean $\pm$ SD over 10 GroupKFold folds), pooled
OOF value, and one-sided one-sample Wilcoxon signed-rank test against the RFE+XGBoost
reference value. Pooled Macro F1 exceeds the fold mean because several tail classes are
concentrated in a few wells: folds whose held-out wells lack a class average a zero F1 for
it, which depresses the per-fold values but not the pooled one.\label{tab:perfold}}
\begin{tabularx}{\textwidth}{Xcccc}
\toprule
\textbf{Metric} & \textbf{Fold mean $\pm$ SD} & \textbf{Pooled OOF} & \textbf{Folds $>$ ref.} & \textbf{Wilcoxon $p$}\\
\midrule
Weighted F1 & $0.7521\pm 0.0477$ & $0.7526$ & 7/10 & $0.080$\\
Macro F1 & $0.4410\pm 0.0842$ & $0.5315$ & --- & ---\\
Boundary F1 & $0.4567\pm 0.0250$ & $0.4567$ & 9/10 & $\mathbf{0.002}$\\
Penalty & $-0.5884\pm 0.0937$ & $-0.5886$ & 6/10 & $0.138$\\
\bottomrule
\end{tabularx}
\end{table}

\begin{figure}[H]
\centering
\includegraphics[width=\linewidth]{figures/fig5_perfold_forest.pdf}
\caption{Per-fold metrics across the 10 GroupKFold folds. Each row is one held-out fold;
filled markers exceed the pooled RFE+XGBoost reference (dashed line) and open markers fall
below. The diamond row gives mean $\pm$ SD. Fold-to-fold spread reflects genuine geological
heterogeneity of the held-out well groups, which is why Section~\ref{sec:rq1} reports
dispersion and significance rather than point estimates alone.}
\label{fig:perfold}
\end{figure}

\subsection{Where the Gains Come From}
\label{sec:rq2}
Because the headline comparison combines several components, we next examine whether the
improvement can be explained solely by a stronger base learner. The analysis uses three forms
of evidence: GeoMERIT's position relative to the published RFE ablation ladder, the consistency
of the gain across tree backbones, and a same-stack contrast that isolates the decision layer.
On the ladder of~\cite{chen2026rfe} (Table~\ref{tab:ladder}), quality-aware weighting and
feature engineering each close part of the gap from the $0.6635$ baseline to $0.7272$, and
GeoMERIT's added layers---the three-backend geometric ensemble, the geological hierarchy, the
tail experts, and Bayes-risk decoding---add a further $+0.0254$ Weighted F1 and $+0.0403$
Penalty. Because the published rungs were not re-run inside our stack, this final increment is
interpreted as the cumulative effect of the added layers rather than as a per-layer
decomposition. The decision layer is isolated causally in the same-stack ablation below, while
the representation and learning layers are stress-tested directly in Sections~\ref{sec:rq3}--\ref{sec:rq4}.

\begin{table}[H]
\caption{Ablation ladder: published RFE rungs (from~\cite{chen2026rfe}) versus GeoMERIT. The
GeoMERIT increment is attributed cumulatively to its ensemble, hierarchy, tail-expert, and
Bayes-risk layers; the decision-layer share is isolated separately in
Table~\ref{tab:main}.\label{tab:ladder}}
\begin{tabularx}{\textwidth}{Xcc}
\toprule
\textbf{Configuration} & \textbf{Weighted F1} & \textbf{Penalty}\\
\midrule
Baseline XGBoost (no FE, no weighting) & 0.6635 & $-0.7804$\\
Baseline $+$ weighting & 0.7230 & $-0.6346$\\
Baseline $+$ feature engineering & 0.7259 & $-0.6327$\\
RFE $+$ XGBoost (FE $+$ weighting) & 0.7272 & $-0.6289$\\
\textbf{GeoMERIT} ($+$ ensemble, hierarchy, tail experts, Bayes-risk decoding) & \textbf{0.7526} & $\mathbf{-0.5886}$\\
\bottomrule
\end{tabularx}
\end{table}

The backbone comparison gives the same conclusion. Across backbones
(Table~\ref{tab:backbones}), RFE-style features improve every single-model variant, consistent
with~\cite{chen2026rfe}, yet each remains below GeoMERIT on all three facets. The residual gap
between the best single model (RFE+LightGBM at $0.7092$) and GeoMERIT ($0.7526$) is therefore
associated with components that do not depend on a single backbone: geometric fusion, the
hierarchy, tail gating, and cost-aware decoding.

\begin{table}[H]
\caption{Tree-backbone reference rungs (from~\cite{chen2026rfe}) versus the GeoMERIT
ensemble. Every single-backbone RFE variant remains below the GeoMERIT operating point on all
three facets, so the gain is not explained by the choice of a single backbone.\label{tab:backbones}}
\begin{tabularx}{\textwidth}{Xccc}
\toprule
\textbf{Model} & \textbf{Weighted F1} & \textbf{Boundary F1} & \textbf{Penalty}\\
\midrule
Random Forest & 0.6918 & 0.3840 & $-0.6776$\\
RFE $+$ Random Forest & 0.7051 & 0.3920 & $-0.6467$\\
CatBoost & 0.6501 & 0.3510 & $-0.7895$\\
RFE $+$ CatBoost & 0.6866 & 0.3725 & $-0.7156$\\
LightGBM & 0.6550 & 0.3768 & $-0.8080$\\
RFE $+$ LightGBM & 0.7092 & 0.4147 & $-0.6691$\\
\textbf{GeoMERIT ensemble} & \textbf{0.7526} & \textbf{0.4567} & $\mathbf{-0.5886}$\\
\bottomrule
\end{tabularx}
\end{table}

The most direct same-stack evidence comes from the decision layer, which can be evaluated while
holding the representation, learning, and posterior layers fixed. The third and fourth rows of
Table~\ref{tab:main} move the decoder from plain Bayes risk to the full logit-adjusted, gated
Bayes decoder. This adds $+0.0080$ Weighted F1, $+0.0277$ Macro F1, $+0.0144$ Boundary F1, and
$+0.0324$ Tail F1 at essentially unchanged Penalty ($-0.5885\to-0.5886$). Because the ensemble,
weighting, and features are identical, this contrast gives a causal estimate of the decision
layer's effect. The decision layer converts posterior quality into rare-class and boundary
performance without increasing geological penalty, which is the intended behavior of the
Bayes-risk formulation in Equation~(\ref{eq:bayes}). The remaining $+0.0174$ Weighted F1
between the base ensemble and the full RFE rung is attributed to the representation, learning,
and posterior layers combined, and is stress-tested in the next two sections.

\subsection{Robustness on Feature-Deficient Wells}
\label{sec:rq3}
A direct test of the representation layer is performance on wells with incomplete curve suites.
We stratify the held-out wells by mean per-row missingness and compare feature-complete against
feature-deficient subsets at several thresholds, annotating each bucket with its size in wells
so that single-well buckets can be interpreted separately from populated ones
(Figure~\ref{fig:robust}). Performance on missingness-heavy wells degrades gradually rather
than collapsing. The $38$-well deficient bucket ($0.1/0.4$ split) attains Weighted F1 $0.7103$
and Boundary F1 $0.4408$, above the full-fleet RFE+XGBoost boundary reference of $0.4100$
despite operating only on incomplete wells. The $8$-well deficient bucket ($0.2/0.5$) attains
Weighted F1 $0.7555$, exceeding the full-fleet reference of $0.7272$. The $0.3/0.6$ deficient
bucket is the only severe drop, but it contains a single well of $11{,}315$ rows and should be
read as a single-well case rather than a systematic threshold effect. These results support the
intended role of the missingness indicators, curve-suite fingerprint, and within-well
normalization: they retain information that remains available when parts of the curve suite are
absent, which is the practical case of a newly drilled well whose tool string differs from the
training fleet.

\begin{figure}[H]
\centering
\includegraphics[width=\linewidth]{figures/fig6_missingness_robustness.pdf}
\caption{Missingness-stratified robustness under held-out-well validation. At each threshold,
wells are split into feature-complete (solid blue) and feature-deficient (dashed orange)
subsets by mean per-row missing rate; annotations give bucket sizes in wells. The gray
dashed line is the pooled full-fleet RFE+XGBoost reference. Deficient buckets with
meaningful population (38 and 8 wells) stay near or above the reference on all facets; the
single-well $0.3/0.6$ deficient bucket explains its own dip.}
\label{fig:robust}
\end{figure}

\subsection{Per-Class Diagnosis, Tail Behavior, and Well Tracks}
\label{sec:rq4}
Because Weighted F1 is dominated by Shale, it can obscure rare-class behavior and
formation-boundary errors. We therefore report per-class out-of-fold precision, recall, and F1
(Table~\ref{tab:perclass}) and interpret them through a three-level diagnosis
(Figure~\ref{fig:classdiag}): \emph{which} classes remain weak (panel~a), \emph{where} their
probability mass is assigned when they are missed (panel~b), and whether the six rare classes
targeted by the tail-expert mechanism are recovered (panel~c).

\begin{table}[H]
\caption{Per-class out-of-fold metrics for the 12 FORCE 2020 lithology classes (pooled over
10 folds). Support is the number of labeled rows. Figure~\ref{fig:classdiag}a visualizes the
F1 column jointly with support; precision and recall are reported here
only.\label{tab:perclass}}
\begin{tabularx}{\textwidth}{Xrccc}
\toprule
\textbf{Class} & \textbf{Support} & \textbf{Precision} & \textbf{Recall} & \textbf{F1}\\
\midrule
Shale & 877{,}043 & 0.833 & 0.912 & 0.870\\
Sandstone & 207{,}704 & 0.768 & 0.723 & 0.745\\
Sandstone/Shale & 180{,}820 & 0.442 & 0.285 & 0.347\\
Limestone & 69{,}498 & 0.636 & 0.659 & 0.647\\
Marl & 41{,}038 & 0.382 & 0.368 & 0.375\\
Tuff & 17{,}431 & 0.663 & 0.742 & 0.700\\
Halite & 14{,}712 & 0.968 & 0.983 & 0.975\\
Chalk & 14{,}043 & 0.815 & 0.228 & 0.357\\
Coal & 4{,}754 & 0.649 & 0.692 & 0.670\\
Dolomite & 2{,}391 & 0.040 & 0.041 & 0.040\\
Anhydrite & 1{,}808 & 0.606 & 0.704 & 0.651\\
Basement & 141 & 0.000 & 0.000 & 0.000\\
\bottomrule
\end{tabularx}
\end{table}

\begin{figure}[H]
\centering
\includegraphics[width=\linewidth]{figures/fig7_class_diagnosis.pdf}
\caption{Class-wise diagnosis of GeoMERIT, three nested levels.
(\textbf{a})~Per-class F1 with marker size proportional to log support (support in
parentheses); classes ordered by geological group with the shared lithology palette.
(\textbf{b})~Row-normalized confusion matrix in the same group order as
Figure~\ref{fig:penalty}; framed blocks show that residual errors occur mainly inside the
clastic and carbonate--marl clusters (values below $0.02$ suppressed).
(\textbf{c})~The six tail classes against the $0.65$ recovery goal: four are recovered
(green); Dolomite and Basement remain failure modes (red).}
\label{fig:classdiag}
\end{figure}

Two structural findings are more informative than the raw ranking of scores. First,
\emph{physically distinctive tail classes are recoverable}. Halite (F1 $0.975$; low GR, very
low density), Tuff ($0.700$; high GR with high density), Coal ($0.670$), and Anhydrite
($0.651$) all exceed $0.65$ despite jointly holding under $3\%$ of labels, indicating that the
tail-expert and weighting layers help recover rare classes with distinctive log responses.
Second, \emph{the remaining errors are geologically adjacent}. The F1 gaps in panel (a) do not
by themselves distinguish rarity from lithological overlap, but panel (b) separates the two by
showing where the missed probability mass is assigned. Sandstone/Shale loses $0.538$ of its
mass to Shale (its end-member); Chalk loses $0.654$ to Limestone; Marl loses $0.408$ to Shale
and $0.178$ to Limestone; Dolomite is absorbed by Shale ($0.437$) and Limestone ($0.311$). All
large off-diagonal entries lie inside the two framed clusters that the penalty matrix
(Figure~\ref{fig:penalty}) treats as lower-cost confusions. The tail mean F1 of $0.5061$ is
roughly $1.6\times$ the $\approx\!0.31$ level of single-model baselines.

To rule out an averaging artifact we inspect two held-out wells chosen to span the two
operating conditions Section~\ref{sec:repr} emphasizes: feature-complete 16/7-6 and
missingness-heavy 35/9-7, both located within the North Sea study area shown in
Figure~\ref{fig:well_locations}.
Figure~\ref{fig:tracks} shows, for each well, five measured curves (GR, RHOB, NPHI, DTC,
RDEP), the per-row suite missing fraction, and the true/predicted/hit lithology strips on a
common depth axis.

\begin{figure}[H]
\centering
\includegraphics[width=\linewidth]{figures/fig8_well_tracks.pdf}
\caption{Held-out single-well prediction tracks.
(\textbf{a})~Feature-complete well 16/7-6 (point accuracy $75.5\%$): five log curves,
per-row suite missing fraction, and true/predicted/hit strips in the shared lithology
palette; dotted lines bracket the 2200--2400~m carbonate interval, where the chalk section
is decoded as Limestone---the depth-localized expression of the
$\text{Chalk}\to\text{Limestone}$ confusion ($0.654$) of Figure~\ref{fig:classdiag}b.
(\textbf{b})~Missingness-heavy well 35/9-7 (point accuracy $87.6\%$): six of twenty suite
curves (RSHA, SGR, SP, DCAL, MUDWEIGHT, RXO) are entirely absent (gray missing-fraction
track), yet the predicted column reproduces the Shale-dominated sequence and its sandstone
interbeds with only thin boundary errors. Hatched bands mark within-curve gaps.}
\label{fig:tracks}
\end{figure}

The feature-complete well illustrates the confusion analysis at depth scale. Prediction quality
is high along the clastic section, while the bracketed 2200--2400~m interval---about 1{,}200
rows of Chalk with thin Marl---is decoded almost entirely as Limestone (interval accuracy
$6.0\%$) and accounts for $48\%$ of the well's errors. The error is therefore consistent with
the carbonate--marl ambiguity isolated in Figure~\ref{fig:classdiag}b, now visible as a
contiguous geological interval whose RHOB/NPHI response lies between the chalk and limestone
training distributions. The missingness-heavy well illustrates the representation claim. With
six curves absent for the entire well, point accuracy is $87.6\%$ and errors remain confined to
thin transitions, consistent with the role of explicit missingness encoding described in
Section~\ref{sec:repr}.

\subsection{Decision-Layer Sensitivity}
\label{sec:rq5}
For a cost-guided decision layer to be practically useful, its selected operating point should
lie in a stable region of the validation surface. We therefore sweep the two decision-stage
parameter groups from Section~\ref{sec:decision}---the logit-adjustment temperature $\tau$ and
the Dolomite gate $(\gamma_\text{dol},\theta_\text{dol})$---while the rest of the pipeline is
held fixed (Figure~\ref{fig:decision}).

\begin{figure}[H]
\centering
\includegraphics[width=\linewidth]{figures/fig9_decision_layer.pdf}
\caption{Decision-layer trade-off under logit adjustment and Dolomite gating.
(\textbf{a})~Sweeping $\tau$ (gates off): the four F1 facets rise as the head-class prior is
flattened while Penalty peaks near $\tau\approx 0.1$ and degrades slowly beyond; the
adopted $\tau=0.2$ (dashed) sits inside the shaded flat region, trading $\le 0.002$ Penalty
for $+0.009$ Macro F1 relative to $\tau=0.1$.
(\textbf{b})~Dolomite-gate grid: each point is one
$(\gamma_\text{dol},\theta_\text{dol})$ setting, colored by $\gamma_\text{dol}$; the red
curve is the Pareto frontier of Dolomite F1 versus Penalty and the star is the adopted
setting $(8.0, 0.05)$. Recall cannot be forced. Pushing Dolomite F1 toward its grid
maximum ($\approx 0.076$) costs up to $0.04$ Penalty.}
\label{fig:decision}
\end{figure}

The two panels characterize complementary aspects of the decision layer. The $\tau$ sweep is
monotone and flat near the adopted value; between $\tau=0.1$ and $\tau=0.3$, Weighted F1 varies
by only $0.0045$ and Penalty by $0.0077$. Thus, the long-tail correction occupies a broad,
auditable region rather than a narrow optimum. The gate grid shows that the targeted Dolomite
mechanism lies \emph{on} the Pareto frontier at the adopted point, but higher Dolomite F1
requires a reduction in geological penalty. This trade-off is consistent with a class whose log
response overlaps its carbonate neighbors (Section~\ref{sec:failures}). The monotone and
geologically interpretable surfaces are a benefit of separating decoding from learning: a
retraining-based cost treatment would have to re-fit the model to expose the same trade-off
curves.

\section{Discussion}
\label{sec:discussion}

\subsection{What the Joint Gains Establish About Robustness}
\label{sec:joint}
The headline result is not any single number but a pattern. Under a protocol that forbids
within-well leakage, the representation, weighting, posterior structure, and decoding were each
adjusted to address a named condition of field well-log data, and the resulting system improved
\emph{overall quality, boundary stability, and geological cost at the same time}. The three
improvements are traceable to different layers. The Weighted F1 gain arises chiefly from the
enriched representation and backbone diversity; the Boundary F1 gain (the largest relative
gain, $+11.4\%$, and the only one individually fold-significant) from window statistics,
gradients, and boundary up-weighting; and the Penalty gain---isolated causally by the
same-stack ablation of Section~\ref{sec:rq2}---from the Bayes-risk decoder. This decomposition
matches the design intent of Table~\ref{tab:layers}. Lithology prediction under complex
well-log conditions becomes more robust when data quality, long-tail structure, boundary
uncertainty, and geological cost are optimized jointly rather than appended as post-processing.

\subsection{Failure Modes Are Geologically Structured}
\label{sec:failures}
The classes GeoMERIT still misses show identifiable geological causes, which suggests targeted
remedies rather than a generic need for more data.

\textbf{Dolomite (F1 $=0.040$)} is the largest visible failure, and its cause is mineralogical
rather than statistical. Clean dolomite (PEF $\approx 3.1$, matrix density
$\approx 2.87$~g/cm$^3$) sits squarely inside the carbonate subspace already occupied by
limestone (PEF $\approx 5$, $\approx 2.71$~g/cm$^3$), chalk, and marl, and the gate sweep of
Section~\ref{sec:rq5} showed that overlap to be strong enough that forcing dolomite recall
actively worsens the penalty. The remedy is therefore not a larger classifier but conditional
structure---a carbonate-subspace expert that operates only within the super-class, retrieval
priors restricted to the same GROUP and FORMATION, or a dedicated dolomite-versus-limestone
discriminator trained only where the two rocks coexist.

\textbf{Basement (F1 $=0.000$)} fails completely, but for a reason that lies in the data rather
than in model capacity. Its 141 rows are concentrated in a few wells where the key curves are
heavily missing, so under GroupKFold several training folds contain no Basement supervision at
all. This suggests that extremely rare, spatially localized classes may be better treated as
open-set or anomaly segments than as an additional softmax output.

\textbf{Sandstone/Shale ($0.347$) and Chalk ($0.357$)} are weak despite ample support---
Sandstone/Shale alone has $180{,}820$ rows---so here the ceiling is set by the labels, not by
rarity. Sandstone/Shale is by definition a transitional rock between its two parents, and
Chalk's combination of high precision ($0.815$) with low recall ($0.228$) describes a model
that commits to chalk only when the evidence is overwhelming and otherwise defers to limestone.
These are label-resolution limits at log scale, and moving past them will take either revisited
labels or sequence-level decoding that reasons about bed packages instead of individual rows.

\subsection{Scope and Limitations}
\label{sec:limits}
Three limitations bound the claims. First, the cross-stack comparison in
Table~\ref{tab:ladder} splices published rungs with our extension, so the four-layer
increment is attributed cumulatively rather than per-layer in that comparison; the
same-stack decision-layer ablation of Section~\ref{sec:rq2} provides the only directly
causal contrast within GeoMERIT, and a full in-stack leave-one-out ablation across all four
layers is the natural next experiment. Second, the Weighted F1 and Penalty gains, while
consistent in pooled OOF and directionally consistent across folds, are not individually
significant at $\alpha=0.05$ even under the one-sided test (Table~\ref{tab:perfold}); we
have deliberately not claimed otherwise. Third, all conclusions are for one basin and one
penalty matrix; the decision layer would inherit any context-dependence of $A$ (e.g., depth-
or formation-conditional costs), which we view as an opportunity---Equation~(\ref{eq:bayes})
accepts any cost matrix without retraining.

\subsection{Why a Tree-Only Stack Remains the Right Substrate Here}
\label{sec:trees}
GeoMERIT instantiates, in a geological setting, the broader observation that tree ensembles
remain highly competitive on tabular
data~\cite{grinsztajn2022trees,shwartzziv2022tabular,mcelfresh2024benchmark}. GBDTs natively
route the sentinel-encoded missingness states, remain insensitive to the
four-orders-of-magnitude scale differences across curves, use axis-aligned decision surfaces
that mirror the cross-plot logic of petrophysical interpretation, and accept categorical
features such as \texttt{curve\_suite\_id} and FORMATION as first-class inputs. None of
this precludes deep follow-ups---well-log foundation models~\cite{qi2025wlfm} could supply
upstream representations that the same weighting, hierarchy, and decoding layers consume,
and agentic workflows~\cite{wang2026geodecider} could invoke GeoMERIT as a calibrated
tool---but in the missingness-heavy, long-tailed, strictly well-grouped regime studied here,
the tree substrate is a modeling choice matched to the data, not a concession to compute.

\section{Conclusions}
\label{sec:conclusion}
We presented GeoMERIT, a missingness-aware and penalty-guided lithology prediction method for
incomplete, heterogeneous, and long-tailed well-log data. GeoMERIT targets stability under
field well-log conditions by encoding missingness, anomalies, and curve-suite identity as an
explicit state; increasing the influence of rare classes and boundary samples during training;
refining posteriors through a tree ensemble, a geological hierarchy, and tail experts; and
selecting labels by minimizing geological risk under the official penalty matrix rather than by
maximizing posterior probability. On the full FORCE 2020 valid-label set under strict 10-fold
GroupKFold validation by well---the protocol that simulates prediction on unseen wells---GeoMERIT
attains Weighted F1 $=0.7526$, Macro F1 $=0.5315$, Boundary F1 $=0.4567$, Penalty $=-0.5886$,
and tail mean F1 $=0.5061$, improving the published RFE+XGBoost reference simultaneously by
$+3.5\%$ (Weighted F1), $+11.4\%$ (Boundary F1, fold-significant at $p=0.002$), and $+6.4\%$
(Penalty). A same-stack ablation isolates the decision layer causally, adding $+0.0080$
Weighted F1, $+0.0277$ Macro F1, $+0.0144$ Boundary F1, and $+0.0324$ Tail F1 at essentially
unchanged Penalty. The main conclusion is structural: the three headline gains correspond to
distinct facets of robustness and are produced by separable layers of the design. As a result,
the framework's components can be adopted independently, and its decision layer can be retuned
to another cost matrix without retraining. Future work follows from the diagnosed failure modes:
Dolomite calls for carbonate-subspace experts and geology-conditioned retrieval priors;
Basement-like classes are better posed as open-set or anomaly segments; semi-Markov segment
decoding would allow bed-scale decisions; and well-log foundation models could supply upstream
representations beneath the same robust decision stack.

\vspace{6pt}

\authorcontributions{Conceptualization, methodology, software, validation, formal analysis,
investigation, data curation, writing---original draft preparation, writing---review and
editing, and visualization: A.A. The author has read and agreed to the published version of
the manuscript.}

\funding{This research received no external funding.}

\institutionalreview{Not applicable.}

\informedconsent{Not applicable.}

\dataavailability{The FORCE 2020 dataset is publicly available at
\url{https://doi.org/10.5281/zenodo.4351156}~\cite{bormann2020force}. The code,
configuration files, fold definitions, out-of-fold predictions, decoding grids, and
plotting scripts that reproduce every number, table, and figure in this paper are described
in Appendix~\ref{app:repro}.}

\acknowledgments{The authors thank the FORCE 2020 organizers for releasing the dataset and
penalty matrix that made this study possible.}

\conflictsofinterest{The author declares no conflicts of interest.}

\abbreviations{Abbreviations}{
The following abbreviations are used in this manuscript:\\

\noindent
\begin{tabular}{@{}ll}
GBDT & Gradient-boosted decision tree\\
RFE  & Robust feature engineering\\
OOF  & Out-of-fold\\
MAD  & Median absolute deviation\\
OvR  & One-vs-rest\\
GR   & Gamma ray\\
RHOB & Bulk density\\
NPHI & Neutron porosity\\
PEF  & Photoelectric factor\\
DTC  & Compressional sonic transit time\\
RDEP & Deep resistivity\\
NPD  & Norwegian Petroleum Directorate\\
\end{tabular}}

\appendixtitles{yes}
\appendixstart
\appendix
\section[\appendixname~\thesection]{Reproducibility Artifacts}
\label{app:repro}
All numbers in this paper are regenerated from saved artifacts. Core library:
\texttt{geomerit/features.py} (representation layer), \texttt{weights.py} (learning layer),
\texttt{models.py} (posterior layer: three backends, geometric fusion, coarse-to-fine, tail
experts), \texttt{decode.py} (decision layer), \texttt{metrics.py} (five metrics),
\texttt{cv.py} (GroupKFold protocol). Result files: \texttt{decode\_report.json} (pooled
metrics and per-class report), \texttt{decode\_report.csv} (per-row OOF predictions),
\texttt{decode\_grid\_results.json} (decision-layer grid), \texttt{perfold\_geomerit.csv}
(per-fold metrics; \texttt{scripts/06\_perfold\_stats.py} reproduces the Wilcoxon tests).
Figures: every panel is generated by \texttt{revision\_v3/scripts/make\_figures.py} from
same-named CSV source data; the two case-well curve extracts derive from the public FORCE
2020 release. Best decoding configuration: stage \texttt{global\_tail}, $\tau=0.2$,
$(\gamma,\theta)_\text{global}=(1.2,0.35)$,
$(\gamma,\theta)_\text{dol}=(8.0,0.05)$.

\begin{adjustwidth}{-\extralength}{0cm}
\reftitle{References}
\bibliography{references}
\end{adjustwidth}

\PublishersNote{}
\end{document}